\documentclass{resume} 
\usepackage{hyperref}
\usepackage{fontawesome5}
\usepackage[dvipsnames]{xcolor}
\hypersetup{colorlinks,linkcolor=red}
\definecolor{mypink1}{rgb}{0.858, 0.188, 0.478}
\definecolor{skillExpert}{HTML}{0072B2}
\definecolor{skillAdvanced}{HTML}{009E73}
\definecolor{skillProficient}{HTML}{E69F00}
\definecolor{skillFamiliar}{HTML}{CC79A7}
\definecolor{skillBeginner}{HTML}{8A6F73}
\newcommand{\skill}[2]{\textcolor{#1}{#2}}
\usepackage[left=0.75in,top=0.6in,right=0.75in,bottom=0.6in]{geometry} % Document margins
\usepackage{mathpazo}
\newcommand{\tab}[1]{\hspace{.2667\textwidth}\rlap{#1}}
\newcommand{\itab}[1]{\hspace{0em}\rlap{#1}}
\name{Yunhai Xiang} 
\address{Mathematician  \(\cdot\) Data Scientist \(\cdot\) Applied Scientist \(\cdot\) Software Engineer}
\address{\small
\href{mailto:yxiang72@uwo.ca}{\faEnvelope\ yxiang72@uwo.ca}\quad
\href{https://yunhaimath.com/}{\faGlobe\ yunhaimath.com}\quad
\href{https://github.com/yunhaixiang}{\faGithub\ yunhaixiang}\quad
\href{https://www.linkedin.com/in/yunhai-xiang/}{\faLinkedin\ yunhai-xiang}}

\begin{document}

%----------------------------------------------------------------------------------------
%	EDUCATION SECTION
%----------------------------------------------------------------------------------------

\begin{rSection}{Education}
\rSubsectionEducation{Western University (University of Western Ontario)}{2024 --- Now}{Ph.D. (in progress) and M.Sc. in Mathematics}{London, ON, Canada}{Advisor: Chris Hall\\Specialization: Number Theory and Arithmetic Geometry}
\smallskip
\rSubsectionEducation{University of Waterloo}{2018 --- 2023}{B.Math. in Pure Mathematics, Honors, with Distinction}{Waterloo, ON, Canada}{}
\end{rSection}
%----------------------------------------------------------------------------------------
%	WORK EXPERIENCE SECTION
%----------------------------------------------------------------------------------------

\begin{rSection}{Industry Experience}
\begin{rSubsection}{ByteDance}{2019}{Co-op Internship}{Beijing, China}
\item Supported data-informed product development for an education technology product
\item Collected and reviewed educational data to support product analysis and application improvement
\end{rSubsection}
\end{rSection}
\vspace{0.25em}
%----------------------------------------------------------------------------------------
%	TECHNICAL STRENGTHS SECTION
%----------------------------------------------------------------------------------------

\begin{rSection}{Skills}
\textit{Colored by level of mastery:
\skill{skillExpert}{Expert},
\skill{skillAdvanced}{Advanced},
\skill{skillProficient}{Proficient},
\skill{skillFamiliar}{Familiar},
\skill{skillBeginner}{Beginner}}
\vspace{-0.4em}

\begin{tabular}{ @{} >{\bfseries}l @{\hspace{6ex}} l }
Languages & \skill{skillExpert}{Mandarin}, \skill{skillExpert}{English}, \skill{skillBeginner}{French}, \skill{skillBeginner}{Japanese}\\
Programming &  \skill{skillAdvanced}{Python}, \skill{skillFamiliar}{Lean}, \skill{skillFamiliar}{C/C++}, \skill{skillBeginner}{C\#}, \skill{skillBeginner}{Rust},  \skill{skillBeginner}{JavaScript/TypeScript}, \skill{skillBeginner}{Haskell}, \skill{skillBeginner}{Scheme} \\
Computing & \skill{skillExpert}{SageMath}, \skill{skillProficient}{NumPy}, \skill{skillProficient}{SciPy}, \skill{skillFamiliar}{Pandas}, \skill{skillFamiliar}{PyTorch}, \skill{skillFamiliar}{XGBoost}, \skill{skillFamiliar}{Sklearn},  \skill{skillFamiliar}{Slurm}\\
DSLs & \skill{skillExpert}{LaTeX}, \skill{skillProficient}{SQL}, \skill{skillProficient}{HTML/CSS}\iffalse, \skill{skillBeginner}{GDScript}\fi\\
Softwares & \skill{skillProficient}{Cinema 4D}, \skill{skillBeginner}{Blender}, \skill{skillBeginner}{Unity}, \skill{skillBeginner}{Godot}\\
Hobbies & \skill{skillProficient}{Magic Tricks}, \skill{skillProficient}{Swimming}, \skill{skillBeginner}{Indie Games}\\
\end{tabular}
\vspace{0.3em}
\end{rSection}
\vspace{0.55em}
\begin{rSection}{Publications \& Preprints}
{\setlength{\leftmarginii}{0pt}
\begin{list}{}{\leftmargin=-0.6em \labelwidth=0pt \labelsep=0pt \itemsep=-2pt}
    \item \textit{QEDBench: Quantifying the Alignment Gap in Automated Evaluation of University-Level Mathematical Proofs}\\
    Joint with Santiago Gonzalez, Alireza Amiri Bavandpour, Peter Ye, Edward Zhang, et al.\\
    \textsc{International Conference on Machine Learning (ICML) 2026}
\end{list}
}
\end{rSection}
\vspace{0.55em}
\begin{rSection}{Awards}
\begin{rSubsection}{Mathematics Undergraduate Research Award}{2022}{University of Waterloo}{}
\item Highly competitive undergrad research award of value \$10000 at University of Waterloo. Worked on research project in summer related to Weyl groups, supervized by Prof. Matthew Satriano.
\end{rSubsection}
\end{rSection}
\vspace{0.25em}

\begin{rSection}{Academic Meetings \& Presentations} \itemsep -3pt
\begin{rSubsection}{Conferences Attendance \& Presentations}{}{Attended academic conferences and delivered research or expository talks}{}
\item $\infty$-Type Café Summer School 2026, Online 
\item FrontierMath Open Problems Workshop,  Summer 2026, UToronto
\item Homotopy Theory, K-theory, and Topological Data Analysis Conference, Summer 2026, Western
\item Using AI to Find Examples Workshop, under thematic program Universal Statistics in Number Theory, Winter 2026, Centre de Recherches Mathématiques, Université de Montréal
\item University of Western Ontario Research in Computer Science (UWORCS), Winter 2026, Western
\item {\em Rigid Analytic Geometry}, Ontario Graduate Mathematics Conference,  Summer 2024, UWaterloo
\item {\em Cohomology of Grassmannians}, Canadian Undergraduate Mathematics Conference, Summer 2022, Université Laval
\end{rSubsection}
\begin{rSubsection}{Seminar Attendance \& Presentations}{}{Attended academic seminars and delivered research or expository talks}{}
    \item {\em Hyperelliptic Jacobians and Hasse–Witt Matrix}, Algebraic Geometry Student Seminar, Western
    \item {\em Trace of Frobenius}, Winter 2026, Graduate Student Seminar, Western
    \item {\em Brauer Groups of Fields Part II}, Winter 2026, Algebraic Geometry Student Seminar, Western
    \item {\em Brauer Groups of Fields Part I}, Winter 2026, Algebraic Geometry Student Seminar, Western
    \item {\em Weyl Group and Root Datum}, Summer 2025, Algebraic Groups Seminar, Western
    \item {\em Hecke Operators}, Summer 2025, Algebraic Groups Seminar, Western
    \item {\em Parabolic and Borel Subgroups}, Summer 2025, Algebraic Groups Seminar, Western
    \item {\em Lie Algebras of Algebraic Groups}, Summer 2025, Algebraic Groups Seminar, Western
    \item {\em Admissible Topology}, Summer 2024, Rigid Analytic Geometry Seminar, Western
    \item {\em Non-Archimedean Fields}, Summer 2024, Rigid Analytic Geometry Seminar, Western
    \item {\em Cones and Affine Toric Varieties}, Summer 2024, Toric Varieties Learning Seminar, Western
    \item {\em An Easy Tour of Galois Cohomology}, Winter 2024, Graduate Student Seminar, Western
    \item {\em Introduction to Derived Categories}, Fall 2022, Algebraic Geometry Working Seminar, UWaterloo
\end{rSubsection}
\end{rSection}
\vspace{0.25em}



\begin{rSection}{Teaching \& Mentoring} \itemsep -3pt
    \begin{rSubsection}{Course Instructorship}{}{}{}
        \item Coming soon. 
    \end{rSubsection}
    \begin{rSubsection}{Teaching Assistantship}{}{Proctored and graded exams, graded homeworks, tutored at Help Centre, and taught tutorial sessions}{}
        \item MATH 3150, \textit{Elementary Number Theory I}, Fall 2024
        \item MATH 2155, \textit{Mathematical Structures}, Fall 2025
        \item CALCULUS 1501, \textit{Calculus II for Mathematical and Physical Sciences}, Winter 2025, Winter 2026, Summer 2026
        \item MATH 1229, \textit{Methods of Matrix Algebra}, Summer 2024, Summer 2025
        \item MATH 1228, \textit{Methods of Finite Mathematics}, Summer 2024, Summer 2025
        \item CALCULUS 1000, \textit{Calculus I}, Winter 2024, Summer 2024, Winter 2025, Fall 2025, Winter 2026
    \end{rSubsection}
    \begin{rSubsection}{Mathematical Outreach Lectures}{}{Gave expository lectures to high school students or undergraduate students on advanced mathematical topics}{}
        \item {\em Aperiodic Tilings}, Summer 2026, Western Summer Math Camp (Grade 11/12)
        \item {\em Surreal Numbers and the Game of Go}, Summer 2025, Western Summer Math Camp (Grade 10)
        \item {\em Mathematical Magic Tricks: Gilbreath's Principle}, Summer 2024, Western Summer Math Camp (Grade 11)
        \item {\em Why is There no Quintic Formula}, Summer 2024, Western Summer Math Camp (Grade 10)
        \item {\em Category Theory Demystified}, Winter 2023, Short Attention Span Math Seminar (Undergraduate)
\end{rSubsection}
    \begin{rSubsection}{Directed Reading Program Mentorship}{}{Mentored undergraduate students on reading projects, culminating in a presentation}{}
        \item Jude McNorgan Lubos, \textit{Introduction to Quadratic Reciprocity}, Winter 2026
    \end{rSubsection}
\end{rSection}
\vspace{0.25em}


\begin{rSection}{Service \& Volunteer Work} \itemsep -3pt
    \begin{rSubsection}{Seminar Organization}{}{Organized or co-organized research or learning seminars}{}
        \item Algebraic Geometry Student Seminar, Summer 2026
        \item Algebraic Geometry Student Seminar, Winter 2026, co-organized with Alex Zwart
        \item Algebraic Geometry Student Seminar, Fall 2024, co-organized with Alex Zwart
        \item Toric Varieties Learning Seminar, Summer 2024, co-organized with Nikolas Rieke
    \end{rSubsection}
\end{rSection}

\begin{rSection}{Selected Software Projects \& Open Source Contributions} \itemsep -3pt
    Coming soon. 
\end{rSection}

\begin{rSection}{Certifications} \itemsep -3pt
    Coming soon. 
\end{rSection}


\end{document}
